% =============================================================================
% ArqeloCV — Canonical Resume Template
% =============================================================================
%
% PURPOSE
% -------
% This is the canonical LaTeX template used by ArqeloCV to generate
% software-engineering resumes.
%
% IMPORTANT RULES FOR AI / TEMPLATE RENDERER
% ------------------------------------------
% 1. Never invent experience, skills, technologies, metrics, dates,
%    certifications, achievements, employers, titles, or education.
%
% 2. Only render sections that contain meaningful content.
%    Empty sections should be omitted entirely.
%
% 3. Section order may be changed according to career stage and relevance.
%
%    Typical experienced candidate:
%      Experience -> Skills -> Projects -> Certifications/Achievements -> Education
%
%    Typical student/new graduate:
%      Education/Experience -> Projects -> Skills -> Certifications/Achievements
%
% 4. Prefer one page for students through mid-level/senior candidates when
%    possible without harming readability.
%
% 5. Experience bullets should:
%      - begin with a strong action verb;
%      - contain one sentence;
%      - preferably occupy 1--2 lines;
%      - explain action, context, and impact where evidence exists;
%      - use STAR / CAR / XYZ principles naturally;
%      - NEVER fabricate quantitative impact.
%
% 6. Professional experience has priority over project experience for
%    candidates who already have meaningful professional experience.
%
% 7. Projects should represent substantial work. Avoid trivial tutorial or
%    coursework projects unless they materially strengthen a student resume.
%
% 8. Certifications should be relevant to the target role or meaningfully
%    strengthen the candidate profile.
%
% 9. Achievements should be genuinely notable. Do not turn routine job
%    responsibilities into "achievements".
%
% 10. Escape LaTeX-sensitive characters in user-generated content:
%       &  %  $  #  _  {  }  ~  ^  \
%
% 11. Preserve factual dates exactly unless the user explicitly corrects them.
%
% =============================================================================

\documentclass[11pt]{article}

% -----------------------------------------------------------------------------
% PAGE LAYOUT
% -----------------------------------------------------------------------------
% North America: use letterpaper.
% India / most other regions: the application may switch this to a4paper.
%
% Keep margins compact but readable. Do not shrink them further merely to
% force excessive content onto one page.
% -----------------------------------------------------------------------------
\usepackage[
  letterpaper,
  top=0.5in,
  bottom=0.5in,
  left=0.5in,
  right=0.5in
]{geometry}

% -----------------------------------------------------------------------------
% TYPOGRAPHY
% -----------------------------------------------------------------------------
\usepackage{XCharter}            % Primary resume font.
\usepackage[T1]{fontenc}         % Better PDF font encoding.
\usepackage[utf8]{inputenc}      % UTF-8 source input.
\usepackage{microtype}           % Subtle typography/spacing improvements.

% -----------------------------------------------------------------------------
% RESUME UTILITIES
% -----------------------------------------------------------------------------
\usepackage{enumitem}            % Compact bullet lists.
\usepackage[hidelinks]{hyperref} % Clickable links without colored boxes.
\usepackage{titlesec}            % Section-heading formatting.

\raggedright
\pagestyle{empty}
\urlstyle{same}
\setlength{\parindent}{0pt}

% -----------------------------------------------------------------------------
% MACHINE-READABLE PDF
% -----------------------------------------------------------------------------
% Important for ATS/text extraction.
% -----------------------------------------------------------------------------
\input{glyphtounicode}
\pdfgentounicode=1

% -----------------------------------------------------------------------------
% SECTION HEADINGS
% -----------------------------------------------------------------------------
\titleformat{\section}
  {\bfseries\large}
  {}
  {0pt}
  {}
  [\vspace{1pt}\titlerule\vspace{-6.5pt}]

% -----------------------------------------------------------------------------
% BULLET LISTS
% -----------------------------------------------------------------------------
\renewcommand\labelitemi{$\vcenter{\hbox{\small$\bullet$}}$}

\setlist[itemize]{
  itemsep=-2pt,
  leftmargin=12pt,
  topsep=7pt
}

% =============================================================================
% RESUME CONTENT
% =============================================================================

\begin{document}

% -----------------------------------------------------------------------------
% NAME
% -----------------------------------------------------------------------------
% Render the candidate's preferred professional name.
% Do not include honorifics unless explicitly requested.
% -----------------------------------------------------------------------------
\centerline{\Huge John Doe}

\vspace{5pt}

% -----------------------------------------------------------------------------
% CONTACT INFORMATION
% -----------------------------------------------------------------------------
% Keep this concise and on one line when possible.
%
% Recommended:
%   Email | Phone | LinkedIn | GitHub | Portfolio
%
% Only render fields actually provided by the candidate.
% Use readable labels instead of extremely long raw URLs.
% -----------------------------------------------------------------------------
\centerline{
  \href{mailto:name@gmail.com}{name@gmail.com}
  \;|\;
  \href{https://linkedin.com/in/name}{linkedin.com/in/name}
  \;|\;
  \href{https://github.com/name}{github.com/name}
  \;|\;
  \href{https://portfolio.com}{portfolio.com}
}

\vspace{-10pt}

% -----------------------------------------------------------------------------
% OPTIONAL SUMMARY
% -----------------------------------------------------------------------------
% Render ONLY when the application determines that a summary adds value,
% e.g. senior candidates, career changers, returning professionals, or
% candidates whose background needs framing.
%
% Maximum: 2--3 concise lines.
%
% Do NOT generate generic statements such as:
% "Passionate software engineer with excellent problem-solving skills."
%
% Example:
%
% \section*{Summary}
% Software engineer with 6+ years of experience building scalable frontend
% platforms with React and TypeScript, with additional experience across
% Node.js services, CI/CD, and cloud-native delivery.
%
% \vspace{-6.5pt}

% -----------------------------------------------------------------------------
% SKILLS
% -----------------------------------------------------------------------------
% Group related skills instead of dumping a flat keyword list.
%
% Use categories appropriate to the candidate, for example:
%   Languages
%   Frontend
%   Backend
%   Databases
%   Cloud / DevOps
%   Tools
%
% Never add a technology solely because it appears in the job description.
% -----------------------------------------------------------------------------
\section*{Skills}

\textbf{Languages:} JavaScript, TypeScript, Java, SQL \\
\textbf{Frontend:} React, Next.js, Redux, Tailwind CSS \\
\textbf{Backend:} Node.js, Spring Boot, REST APIs \\
\textbf{Data \& DevOps:} PostgreSQL, Docker, GitHub Actions, Azure

\vspace{-6.5pt}

% -----------------------------------------------------------------------------
% EXPERIENCE
% -----------------------------------------------------------------------------
% For candidates with meaningful professional experience, this should usually
% be the strongest section of the resume.
%
% Format:
%   Job Title, Company -- Location                     Mon YYYY -- Present
%
% Bullet guidance:
%   - Prefer 3--5 strong bullets for recent/relevant roles.
%   - Older or less relevant roles may use fewer bullets.
%   - Start with action verbs.
%   - Describe accomplishments and technical contribution.
%   - Use metrics only when supported by candidate evidence.
%   - Do not invent scope, ownership, leadership, or impact.
% -----------------------------------------------------------------------------
\section*{Experience}

\textbf{Software Engineer,} {Company Name} -- City, ST
\hfill June 2022 -- Present \\

\vspace{-9pt}

\begin{itemize}
  \item Built and maintained production features using React and TypeScript, improving product usability across key customer workflows.
  \item Developed reusable frontend components and shared application patterns to reduce duplication across product surfaces.
  \item Collaborated with product, design, backend, and QA teams to deliver features from requirements through production release.
\end{itemize}

\textbf{Software Engineer,} {Previous Company} -- City, ST
\hfill Jan 2021 -- May 2022 \\

\vspace{-9pt}

\begin{itemize}
  \item Developed customer-facing application features while working within an established engineering and release workflow.
  \item Improved maintainability by refactoring reusable components and reducing duplicated frontend logic.
  \item Investigated production issues and implemented fixes across frontend application flows.
\end{itemize}

\vspace{-18.5pt}

% -----------------------------------------------------------------------------
% PROJECTS
% -----------------------------------------------------------------------------
% Project priority depends on career stage.
%
% Students / new graduates:
%   Projects may be a primary source of technical evidence.
%
% Experienced candidates:
%   Include only projects that add meaningful technical evidence not already
%   demonstrated by professional experience.
%
% Good projects generally demonstrate:
%   - non-trivial engineering work;
%   - practical usage;
%   - meaningful architecture/technical decisions;
%   - active development or real users;
%   - technologies relevant to the candidate's target roles.
%
% Do not describe mandatory tutorials as major independent projects.
% -----------------------------------------------------------------------------
\section*{Projects}

\textbf{Project Title}
\hfill
\href{https://github.com/name/project}{github.com/name/project} \\

\vspace{-9pt}

\begin{itemize}
  \item Built a full-stack application using Next.js, TypeScript, PostgreSQL, and Docker to solve a specific user problem.
  \item Designed the application architecture and implemented core workflows, API integration, persistence, and deployment.
\end{itemize}

\textbf{Project Title}
\hfill
\href{https://project.com}{project.com} \\

\vspace{-9pt}

\begin{itemize}
  \item Developed and maintained a practical software tool used to automate or simplify a recurring engineering workflow.
  \item Implemented relevant technical features while maintaining documentation, source control, and release history.
\end{itemize}

\vspace{-18.5pt}

% -----------------------------------------------------------------------------
% CERTIFICATIONS — OPTIONAL
% -----------------------------------------------------------------------------
% Render this section only when certifications meaningfully strengthen the
% candidate's resume.
%
% Preferred format:
%   Certification Name -- Issuer                         Mon YYYY
%
% A credential link may be attached to the certification name when available.
%
% Do not list:
%   - expired certifications without context;
%   - trivial course-completion certificates;
%   - certificates unrelated to the candidate's target roles merely to fill space.
%
% If space is limited, relevant experience/projects should usually take priority.
% -----------------------------------------------------------------------------
\section*{Certifications}

\textbf{Certification Name} -- Issuing Organization
\hfill June 2024 \\

\textbf{\href{https://credential.example.com}{Certification Name}}
-- Issuing Organization
\hfill Jan 2023

\vspace{-6.5pt}

% -----------------------------------------------------------------------------
% ACHIEVEMENTS — OPTIONAL
% -----------------------------------------------------------------------------
% Render only genuine, notable achievements.
%
% Examples:
%   - competitive programming results;
%   - hackathon wins;
%   - meaningful open-source recognition;
%   - patents;
%   - significant scholarships/academic honors;
%   - major company awards;
%   - recognized technical accomplishments.
%
% Keep entries short.
% Prefer one line per achievement where possible.
%
% Do NOT duplicate achievements already communicated strongly within
% Experience or Projects.
% -----------------------------------------------------------------------------
\section*{Achievements}

\begin{itemize}
  \item Achieved a notable engineering, academic, open-source, or competitive programming accomplishment with relevant context.
  \item Received a meaningful technical award or recognition from an identifiable organization.
\end{itemize}

\vspace{-18.5pt}

% -----------------------------------------------------------------------------
% EDUCATION
% -----------------------------------------------------------------------------
% Usually keep this compact for experienced candidates.
%
% Students/new graduates may include:
%   - expected graduation date;
%   - GPA only when strong/relevant;
%   - selected honors;
%   - relevant coursework only when it adds meaningful signal.
%
% Experienced engineers normally do not need coursework.
% -----------------------------------------------------------------------------
\section*{Education}

\textbf{School Name} -- PhD in Aerospace Engineering
\hfill June 2010 \\

\textbf{School Name} -- MS in Aerospace Engineering
\hfill June 2006 \\

\textbf{School Name} -- BS in Aerospace Engineering
\hfill Apr 2004

\end{document}
